\section{Matrizen und Abbildungen}
In diesem Kapitel werden wir sehen, wie man lineare Abbildungen durch sogenannte Matrizen darstellen kann. Ausserdem werden wir lernen, wie man mit Matrizen rechnet. Doch zuerst müssen wir klären, was eine Matrix eigentlich ist:
\begin{bsp}
{}{}
Sei 
\[
A = \begin{pmatrix} 2 & -4 \\ 3 & 5 \end{pmatrix}
\]
Eine solche Zahlentabelle mit zwei Zeilen und zwei Spalten nennt man eine \( 2 \times 2 \)-Matrix. Die Zahl in der 2. Zeile und der 1. Spalte nennt man \( a_{21} = 3 \).

Diese Matrix kann dazu verwendet werden, eine Abbildung von \( \mathbb{R}^2 \) nach \( \mathbb{R}^2 \) zu beschreiben. Dazu definiert man das Produkt einer Matrix mit einem Vektor:

\[
A \cdot \vec{v} =
\begin{pmatrix} 2 & -4 \\ 3 & 5 \end{pmatrix}
\cdot
\begin{pmatrix} x \\ y \end{pmatrix}
=
\begin{pmatrix} 2x - 4y \\ 3x + 5y \end{pmatrix}
\]

Das Bild des Vektors \( \vec{v} = \begin{pmatrix} 1 \\ 2 \end{pmatrix} \) unter der Abbildung \( f(\vec{v}) = A \cdot \vec{v} \) ist dann:

\[
f \begin{pmatrix} 1 \\ 2 \end{pmatrix} =
\begin{pmatrix} 2 & -4 \\ 3 & 5 \end{pmatrix}
\cdot
\begin{pmatrix} 1 \\ 2 \end{pmatrix}
=
\begin{pmatrix} 2 \cdot 1 - 4 \cdot 2 \\ 3 \cdot 1 + 5 \cdot 2 \end{pmatrix}
=
\begin{pmatrix} -6 \\ 13 \end{pmatrix}
\]
\end{bsp}
\begin{definition}
    {Matrizen}{}
    Eine Matrix mit 2 Zeilen und 2 Spalten nennt man eine \( 2 \times 2 \)-Matrix:

\[
A =
\begin{pmatrix} 
a_{11} & a_{12} \\ 
a_{21} & a_{22} 
\end{pmatrix}
\]

Analog hat eine \( 3 \times 3 \)-Matrix die Form:

\[
A =
\begin{pmatrix} 
a_{11} & a_{12} & a_{13} \\ 
a_{21} & a_{22} & a_{23} \\ 
a_{31} & a_{32} & a_{33} 
\end{pmatrix}
\]

\end{definition}
\begin{definition}
    {Matrix-Vektor-Produkt}{}
    Eine \( 2 \times 2 \)-Matrix \( A \) kann man wie folgt mit einem Vektor in \( \mathbb{R}^2 \) multiplizieren:

\[
A \cdot \vec{v} =
\begin{pmatrix} 
a_{11} & a_{12} \\ 
a_{21} & a_{22} 
\end{pmatrix}
\cdot
\begin{pmatrix} x \\ y \end{pmatrix}
=
\begin{pmatrix} 
a_{11} \cdot x + a_{12} \cdot y \\ 
a_{21} \cdot x + a_{22} \cdot y 
\end{pmatrix}
\]

Analog für eine \( 3 \times 3 \)-Matrix und einen Vektor in \( \mathbb{R}^3 \):

\[
A \cdot \vec{v} =
\begin{pmatrix} 
a_{11} & a_{12} & a_{13} \\ 
a_{21} & a_{22} & a_{23} \\ 
a_{31} & a_{32} & a_{33} 
\end{pmatrix}
\cdot
\begin{pmatrix} x \\ y \\ z \end{pmatrix}
=
\begin{pmatrix} 
a_{11} \cdot x + a_{12} \cdot y + a_{13} \cdot z \\ 
a_{21} \cdot x + a_{22} \cdot y + a_{23} \cdot z \\ 
a_{31} \cdot x + a_{32} \cdot y + a_{33} \cdot z 
\end{pmatrix}
\]

\end{definition}
\clearpage
\begin{ueb}
    {}{}
    \begin{enumerate}
        \item Berechnen Sie:

        \[
        \begin{pmatrix} 1 & 2 \\ 4 & 1 \end{pmatrix}
        \cdot
        \begin{pmatrix} 3 \\ 5 \end{pmatrix}
        \]
        \item Berechnen Sie:

        \[
        \begin{pmatrix} 1 & 0 & -7 \\ -3 & 4 & 1 \\ 2 & 1 & 0 \end{pmatrix}
        \cdot
        \begin{pmatrix} -1 \\ 2 \\ 1 \end{pmatrix}
        \]
    \end{enumerate}

    Kontrolle: $\begin{pmatrix} 13 \\ 17 \end{pmatrix}$ und $\begin{pmatrix} -8 \\ 12 \\ 0 \end{pmatrix}$
\end{ueb}

In \Cref{uebng:lineareAbbildung} haben wir gesehen, dass eine lineare Abbildung durch ihre Wirkung auf die Basisvektoren vollständig definiert wird. Betrachten wir dazu eine lineare Abbildung 
\( f: \mathbb{R}^2 \to \mathbb{R}^2 \) mit 

% \[
% f \begin{pmatrix} 1 \\ 0 \end{pmatrix} = \begin{pmatrix} a_{11} \\ a_{21} \end{pmatrix}, 
% \quad f \begin{pmatrix} 0 \\ 1 \end{pmatrix} = \begin{pmatrix} a_{12} \\ a_{22} \end{pmatrix}
% \]

\[
f \begin{pmatrix} 1 \\ 0 \end{pmatrix} = \begin{pmatrix} a \\ b \end{pmatrix}, 
\quad f \begin{pmatrix} 0 \\ 1 \end{pmatrix} = \begin{pmatrix} c \\ d \end{pmatrix}
\]



Dann gilt wegen der Linearität:



% \begin{align*}
%     f \begin{pmatrix} x \\ y \end{pmatrix}     &= f \left( x \cdot \begin{pmatrix} 1 \\ 0 \end{pmatrix} + y \cdot \begin{pmatrix} 0 \\ 1 \end{pmatrix} \right) \\
%     &= x \cdot f \begin{pmatrix} 1 \\ 0 \end{pmatrix} + y \cdot f \begin{pmatrix} 0 \\ 1 \end{pmatrix} \\
%     &= x \cdot \begin{pmatrix} a_{11} \\ a_{21} \end{pmatrix} + y \cdot \begin{pmatrix} a_{12} \\ a_{22} \end{pmatrix} \\
%     &= \begin{pmatrix} a_{11} \cdot x + a_{12} \cdot y \\ a_{21} \cdot x + a_{22} \cdot y \end{pmatrix} \\
%     &= \begin{pmatrix} a_{11} & a_{12} \\ a_{21} & a_{22} \end{pmatrix} \cdot \begin{pmatrix} x \\ y \end{pmatrix}
% \end{align*}

\begin{align*}
    f \begin{pmatrix} x \\ y \end{pmatrix}     &= f \left( x \cdot \begin{pmatrix} 1 \\ 0 \end{pmatrix} + y \cdot \begin{pmatrix} 0 \\ 1 \end{pmatrix} \right) \\
    &= x \cdot f \begin{pmatrix} 1 \\ 0 \end{pmatrix} + y \cdot f \begin{pmatrix} 0 \\ 1 \end{pmatrix} \\
    &= x \cdot \begin{pmatrix} a \\ b \end{pmatrix} + y \cdot \begin{pmatrix} c \\ d \end{pmatrix} \\
    &= \begin{pmatrix} a \cdot x + c \cdot y \\ b \cdot x + d \cdot y \end{pmatrix} \\
    &= \begin{pmatrix} a & c\\ b & d \end{pmatrix} \cdot \begin{pmatrix} x \\ y \end{pmatrix}
\end{align*}



Dies zeigt, dass jede lineare Abbildung von \( \mathbb{R}^2 \) nach \( \mathbb{R}^2 \) durch eine \( 2 \times 2 \)-Matrix dargestellt werden kann.
\clearpage
\begin{satz}
{Lineare Abbildungen und Matrizen}{}
Eine Abbildung \( f: \mathbb{R}^2 \to \mathbb{R}^2 \) ist genau dann linear, wenn es eine \( 2 \times 2 \)-Matrix \( A \) gibt mit:

\[
f(\vec{v}) = A \cdot \vec{v}, \quad \text{für alle } \vec{v} \in \mathbb{R}^2
\]

wobei die Spalten der Matrix \( A \) die Bilder der Basisvektoren \( \vec{e}_1 = \begin{pmatrix} 1 \\ 0 \end{pmatrix} \) und \( \vec{e}_2 = \begin{pmatrix} 0 \\ 1 \end{pmatrix} \) sind.

% \[
% \vec{e}_1 = \begin{pmatrix} 1 \\ 0 \end{pmatrix}, \quad \vec{e}_2 = \begin{pmatrix} 0 \\ 1 \end{pmatrix}
% \]

Das Gleiche gilt allgemein für Abbildungen von \( \mathbb{R}^m \to \mathbb{R}^n \), wobei es sich dann um eine \( n \times m \)-Matrix handelt.

% Man nennt \( A \) die Matrix der linearen Abbildung \( f \).


\end{satz}

Dieser Satz ist nützlich, wenn es darum geht, die Matrix einer gegebenen Abbildung zu ermitteln. Im Grunde genommen reicht es, wenn man die Bilder der beiden Basisvektoren bestimmt. 

Wie das genau funktioniert, wird im folgenden Beispiel gezeigt.

\begin{bsp}
    {Matrix einer Drehung}{}
    Die Drehung um den Ursprung um \( 30^\circ \) ist eine lineare Abbildung. Wie lautet die dazugehörige Matrix? Und wohin kommt der Punkt \( P(2, 4) \) zu liegen, wenn man ihn dreht?
    
    \vspace*{12 cm}


%     Betrachten Sie die Skizze rechts.  
% Der Vektor \( \vec{e}_1 = \begin{pmatrix} 1 \\ 0 \end{pmatrix} \) wird abgebildet auf  

% \[
% \vec{e}_1' = \begin{pmatrix} \cos 30^\circ \\ \sin 30^\circ \end{pmatrix} = \begin{pmatrix} \frac{\sqrt{3}}{2} \\ \frac{1}{2} \end{pmatrix}.
% \]

% Der Vektor \( \vec{e}_2 = \begin{pmatrix} 0 \\ 1 \end{pmatrix} \) wird abgebildet auf  

% \[
% \vec{e}_2' = \begin{pmatrix} \cos (90^\circ + 30^\circ) \\ \sin (90^\circ + 30^\circ) \end{pmatrix} = \begin{pmatrix} -\frac{1}{2} \\ \frac{\sqrt{3}}{2} \end{pmatrix}.
% \]

% Damit lautet die Matrix der Drehung:

% \[
% A = \begin{pmatrix} \frac{\sqrt{3}}{2} & -\frac{1}{2} \\ \frac{1}{2} & \frac{\sqrt{3}}{2} \end{pmatrix}.
% \]

% Das Bild von \( P \) unter der Drehung ist:

% \[
% A \cdot \begin{pmatrix} 2 \\ 4 \end{pmatrix} =
% \begin{pmatrix} \frac{\sqrt{3}}{2} & -\frac{1}{2} \\ \frac{1}{2} & \frac{\sqrt{3}}{2} \end{pmatrix}
% \cdot
% \begin{pmatrix} 2 \\ 4 \end{pmatrix}
% =
% \begin{pmatrix} \sqrt{3} - 2 \\ 1 + 2\sqrt{3} \end{pmatrix}.
% \]

\end{bsp}
\textbf{Bemerkung:} Wenn Sie mit dem Computer ein Bild um \( 30^\circ \) drehen, wird im Hintergrund genau diese Rechnung durchgeführt. Die Grafikkarten sind im Grunde Bauteile, die effizient mit Matrizen rechnen!
\clearpage
\begin{satz}
    {Drehmatrizen}{}
    Die Matrix einer Drehung um den Ursprung um den Winkel \( \varphi \) im Gegenuhrzeigersinn lautet:

    \vspace*{2cm}

\end{satz}

\begin{ueb}
    {}{}
    \begin{enumerate}
        \item Berechnen Sie die Bilder von \( \begin{pmatrix} 2 \\ 5 \end{pmatrix} \) und \( \begin{pmatrix} -1 \\ 3 \end{pmatrix} \), wenn die lineare Abbildung durch die folgende Matrix gegeben ist:
        \begin{enumerate}
            \item \( \begin{pmatrix} 1 & 2 \\ -1 & 1 \end{pmatrix} \)
            \item \( \begin{pmatrix} 2 & 3 \\ 1 & -2 \end{pmatrix} \)
            \item \( \begin{pmatrix} 4 & -2 \\ -2 & 3 \end{pmatrix} \)
            \item \( \begin{pmatrix} 1 & -2 \\ 0 & 2 \end{pmatrix} \)
        \end{enumerate}
        
        \item Berechnen Sie die Bilder von \( \begin{pmatrix} 1 \\ 2 \\ 3 \end{pmatrix} \) und \( \begin{pmatrix} 1 \\ -4 \\ 2 \end{pmatrix} \) unter der linearen Abbildung:
        
        \[
        f(\vec{v}) = 
        \begin{pmatrix} 1 & 0 & -1 \\ 2 & 1 & 0 \\ 3 & -3 & 1 \end{pmatrix}
        \cdot \vec{v}
        \]
    
        \item Bestimmen Sie die Matrizen der folgenden Abbildungen in \( \mathbb{R}^2 \):
        \begin{enumerate}
            \item Spiegelung an der $x$-Achse.
            \item Spiegelung an der Geraden \( y = x \).
            \item Streckung mit Zentrum im Ursprung um den Faktor \( k=3 \).
            \item Drehung um \( 270^\circ \) um den Ursprung.
            \item Projektion in Richtung \( \begin{pmatrix} -1 \\ 1 \end{pmatrix} \) auf die $x$-Achse.
            \item Projektion parallel zur y-Achse auf die Gerade \( y = x \).
            \item Identität, d.h. die Abbildung mit \( f(\vec{v}) = \vec{v} \) für alle \( \vec{v} \in \mathbb{R}^2 \).
            \item Drehung um \( 45^\circ \) um den Ursprung.
        \end{enumerate}
    
        \item Bestimmen Sie die Matrizen der folgenden Abbildungen in \( \mathbb{R}^3 \):
        \begin{enumerate}
            \item Spiegelung an der $xy$-Ebene.
            \item Drehung um \( 90^\circ \) um die $z$-Achse.
            \item Drehung um \( 60^\circ \) um die $x$-Achse.
        \end{enumerate}
    \end{enumerate}
    

\end{ueb}

\begin{ueb}
    {Programmieraufgabe}{}
    Der folgende Python-Code definiert zwei Klassen für Vektoren in $\R^2$ und $2\times2$-Matrizen.
    
    \lstinputlisting{Programme/MatrixKlasse-tmplt.py}

    \begin{enumerate}
        \item Ergänzen Sie beide Klassen um die Funktion \lstinline|__str__(self)|, damit Matrix- und Vektorobjekte bei Verwendung von \lstinline|print()| schön ausgegeben werden.
        \item Ergänzen Sie die Klasse der Matrizen um eine Methode \lstinline|mal_vec(self, v)|, die als Returnwert ein Vektorobjekt zurückgibt, das dem Produkt aus Matrix mal Vektor entspricht.
    \end{enumerate}


\end{ueb}
\clearpage
\emph{Lösungen}
\begin{enumerate}
    \item Bilder bestimmen...
    \begin{enumerate}
        \item \( \begin{pmatrix} 12 \\ 3 \end{pmatrix} \) und \( \begin{pmatrix} 5 \\ 4 \end{pmatrix} \)
        \item \( \begin{pmatrix} 19 \\ -8 \end{pmatrix} \) und \( \begin{pmatrix} 7 \\ -7 \end{pmatrix} \)
        \item \( \begin{pmatrix} -2 \\ 11 \end{pmatrix} \) und \( \begin{pmatrix} -10 \\ 11 \end{pmatrix} \)
        \item \( \begin{pmatrix} -8 \\ 10 \end{pmatrix} \) und \( \begin{pmatrix} -7 \\ 6 \end{pmatrix} \)
       
    \end{enumerate}
     \item \( \begin{pmatrix} -2 \\ 4 \\ 0 \end{pmatrix} \) und \( \begin{pmatrix} -1 \\ -2 \\ 17 \end{pmatrix} \)q

    \item \( 2 \times 2 \) Matrizen bestimmen:
    \begin{enumerate}
        \item \( \begin{pmatrix} 1 & 0 \\ 0 & -1 \end{pmatrix} \)
        \item \( \begin{pmatrix} 0 & 1 \\ 1 & 0 \end{pmatrix} \)
        \item \( \begin{pmatrix} 3 & 0 \\ 0 & 3 \end{pmatrix} \)
        \item \( \begin{pmatrix} 0 & 1 \\ -1 & 0 \end{pmatrix} \)
        \item \( \begin{pmatrix} 1 & 1 \\ 0 & 0 \end{pmatrix} \)
        \item \( \begin{pmatrix} 1 & 0 \\ 1 & 0 \end{pmatrix} \)
        \item \( \begin{pmatrix} 1 & 0 \\ 0 & 1 \end{pmatrix} \)
        \item \( \begin{pmatrix} \frac{\sqrt{2}}{2} & -\frac{\sqrt{2}}{2} \\ \frac{\sqrt{2}}{2} & \frac{\sqrt{2}}{2} \end{pmatrix} \)
    \end{enumerate}

    \item \( 3 \times 3 \) Matrizen bestimmen:
    \begin{enumerate}
        \item \( \begin{pmatrix} 1 & 0 & 0 \\ 0 & 1 & 0 \\ 0 & 0 & -1 \end{pmatrix} \)
        \item \( \begin{pmatrix} 0 & -1 & 0 \\ 1 & 0 & 0 \\ 0 & 0 & 1 \end{pmatrix} \)
        \item \( \begin{pmatrix} 1 & 0 & 0 \\ 0 & \frac{1}{2} & -\frac{\sqrt{3}}{2} \\ 0 & \frac{\sqrt{3}}{2} & \frac{1}{2} \end{pmatrix} \)
    \end{enumerate}
\end{enumerate}

